\section{Las bets de las grandes empresas tecnológicas y el emprendimiento: por qué los investigadores fundan empresas}

Su Yu (苏煜) considera que antes las bets de cada empresa en torno a los Agents eran muy distintas entre sí, pero que ahora tienden a coincidir, porque el camino de Anthropic se convirtió en el modelo a seguir para la industria. Tanto Anthropic como OpenAI están convergiendo hacia la dirección de la productivity; Google cuenta con modelos sólidos y con una posición privilegiada en su ecosistema, aunque la repercusión de sus productos no necesariamente está a la altura; xAI, entre otras, también presta atención al computer user agent o al general digital agent.




\subsection{La estructura de recursos de la universidad, las empresas y el emprendimiento}

Su Yu explica por qué pasó de la universidad al emprendimiento: la universidad es adecuada para las weird ideas, la proof of concept y los marcos conceptuales; pero cuando los Agents entran en una etapa intensiva en recursos, en la que se necesitan GPU, API, un equipo sólido y una ejecución de ingeniería rápida, la estructura de recursos de la universidad ya no se ajusta. Emprender no es simplemente una cuestión de dinero, sino una forma de seguir investigando; lo que cambió es la forma que adopta la investigación.

\noindent\begin{tabularx}{\textwidth}{p{0.18\textwidth}p{0.36\textwidth}X}
\toprule
Ámbito & Para qué es adecuado & Para qué no es adecuado \\
\midrule
Universidad & Marcos conceptuales, validación de bajo costo, weird ideas de largo plazo & Ingeniería a gran escala, iteración rápida de productos, experimentos que exigen muchos recursos. \\
Grandes empresas tecnológicas & Entrenamiento de modelos grandes, infraestructura, integración del ecosistema & La exploración libre en múltiples direcciones puede verse limitada por los objetivos de la organización. \\
Startup & Ensayo y error rápido, enfoque en el producto, construcción de sistemas completos en torno a los Agents & Fuerte presión de recursos; es necesario encontrar un wedge claro y una vía comercial. \\
\bottomrule
\end{tabularx}

\subsection{El valor del conceptual framework}

Al final de la entrevista, Su Yu dice que le gusta construir conceptual frameworks: no es que tenga una memoria especialmente buena ni que reaccione con especial rapidez, sino que puede aprender muchas cosas, enlazarlas y ver las conexiones entre ellas. Eso es también lo más valioso de esta revisión de conjunto: coloca la historia de los Agents, el lenguaje, las herramientas, el coding, el continual learning, el world model y la difusión en la industria dentro de un marco unificado.

\begin{importantbox}{Por qué vale la pena convertir este episodio en una nota de alto nivel}
Muchos podcasts son densos en opiniones pero de estructura laxa; este episodio es justo lo contrario, porque en sí mismo construye un marco conceptual. Lo importante al organizarlo no es repetir cada ejemplo, sino conservar este marco: la definición de Agent, la historia técnica, la pila del Language Agent, la disolución de fronteras, el cuello de botella del aprendizaje continuo y las bets de la industria.
\end{importantbox}

\subsection{Resumen de la sección}

Los Agents pasaron de la proof of concept a una etapa intensiva en recursos. La división del trabajo entre la universidad, las grandes empresas tecnológicas y las startups está cambiando; quienes sean capaces de construir marcos conceptuales y, a la vez, movilizar recursos de ingeniería tendrán más facilidad para impulsar la siguiente etapa de los Agents.
